% -----------------------------------------------------------------------------
% Results update: E3 status while final run is in progress
% -----------------------------------------------------------------------------
\section{E3: Factual-signal intervention and durability}
\label{sec:results_e3_intervention}

E3 is the causal core of the thesis. Whereas E1 and E2 are observational, E3 varies the checkpoint at which a synthetic factual signal is injected and then measures how much of that signal is retained after matched continuation and adversarial degradation. The design is intentionally non-circular: the candidate window is located using weight-space indicators from E1, while the outcome is behavioural durability of a held-out factual association signal.

The final E3 factual run uses full-parameter fine-tuning from Pythia-160M checkpoints. The injected signal consists of synthetic factual associations between fictional entities, relations, and values. Each taught fact is presented through training templates, while held-out probes query the same taught association through an unseen question template. Control facts use disjoint entities and values. This design avoids the earlier pilot failure in which training and probe facts were inadvertently disjoint: in the final design, probes test generalisation across templates for taught facts, not extrapolation to unrelated facts.

Each cell loads one checkpoint and one signal seed, scores the base model, injects the signal, measures uptake, and then branches from the post-injection model into two durability tests. The first branch performs a fixed amount of continuation training on a frozen Pile-validation-derived corpus reused identically across stages and seeds. The second branch performs a conflicting-fact degradation attack, in which the same entity--relation pairs are fine-tuned with wrong values. Poison budgets are nested and measured in contradicted facts, so the degradation curve is interpretable as the amount of conflicting information needed to weaken or overwrite the injected association.

Before the final thesis sweep, a calibration run verified that the corrected signal construction produces positive uptake and that both the clean-continuation and degradation branches execute correctly. The calibration also showed that raw uptake is stage-dependent, with later checkpoints learning the synthetic facts more strongly under the fixed injection budget. Consequently, the final analysis reports uptake explicitly and treats uptake-normalised durability as the central comparison. The relevant question is not which stage achieves the highest raw post-injection score, but which stage retains the largest fraction of the injected signal after matched continuation and adversarial pressure.

The final E3 result will therefore be reported as three linked curves: uptake by injection stage, normalised clean retention by injection stage, and normalised degradation-resistance by injection stage. Evidence for a sensitive period would require a durability pattern that is better described by a boundary-aligned break near the E1 window than by a smooth monotone ``earlier is better'' trend. If the durability curve declines smoothly with stage, the result will be interpreted as an honest early-matters effect rather than as evidence for a distinct sensitive period.
